
\documentclass[a4paper,10pt]{report}

%\usepackage[landscape]{geometry}
\usepackage[cm-default]{fontspec}
\setromanfont{FreeSerif}
\setsansfont{FreeSans}
\setmonofont{FreeMono}
\usepackage[utf8x]{inputenc}
\usepackage{fontspec}
\usepackage{xunicode}
\usepackage{xltxtra}
\usepackage{xgreek}
\usepackage{amsmath}
\usepackage{unicode-math}
\usepackage{ulem}
\usepackage{color}
\usepackage{verbatim}
\usepackage{nopageno}
\usepackage{graphicx}
\usepackage{textpos}
\setlength{\TPHorizModule}{1cm}
\setlength{\TPVertModule}{1cm}
%\usepackage[colorgrid,texcoord]{eso-pic}
\usepackage[outline]{contour}
\usepackage{wrapfig}
\usepackage{url}
\usepackage{color}
\usepackage{tikz}
\usepackage{fancybox,fancyhdr}
\usepackage{subfigure}
\usepackage{pstricks}
\usepackage{epsfig}
\usepackage{multicol}
\usepackage{listings}
\usepackage{enumerate}
\usepackage{hyperref}
\hypersetup{
  bookmarks=true,
  bookmarksopen=true,
  pdfborder=false,
  pdfpagemode=UseNone,
  raiselinks=true,
  pdfhighlight={/P},
  colorlinks,
  citecolor=black,
  filecolor=black,
  linkcolor=black,
  urlcolor=black
}


\usepackage{fontspec}
\usepackage{xunicode}
\usepackage{xltxtra}


% Margins
%\setlength{\textwidth}{24cm}

%\setlength{\voffset}{0in}
%\setlength{\textheight}{6.8in}
% Colors
\definecolor{cbrown}{rgb}{0.49,0.24,0.07}
\definecolor{cmpez}{rgb}{0.92,0.88,0.79}
\definecolor{cmpez2}{rgb}{0.62,0.33,0.34}
\definecolor{cwhite}{rgb}{1,1,1}
\definecolor{cblack}{rgb}{0,0,0}
\definecolor{cred}{rgb}{0.9,0.15,0.15}
\definecolor{clblue}{rgb}{0.57,0.85,0.97}
\definecolor{corange}{rgb}{0.86,0.49,0.18}
\definecolor{clorange}{rgb}{0.95,0.84,0.61}
\definecolor{cyellow}{rgb}{1,0.95,0.16}
\definecolor{cgreen}{rgb}{0.69,0.8,0.31}
\definecolor{clbrown}{rgb}{0.78,0.67,0.43}
\definecolor{cmagenta}{rgb}{0.79,0.05,0.54}
\definecolor{cgray}{rgb}{0.76,0.73,0.63}
\definecolor{bluesite}{rgb}{0.05,0.43,0.67}
\definecolor{redsite}{rgb}{0.90,0.15,0.15}
\definecolor{backsite}{rgb}{0.07,0.07,0.07}
% ----------- from Costas -----------------------------

%---------------   Main Fonts    -----------------%   

% \setmainfont[Mapping=tex-text]{Calibri}
% \setmainfont[Mapping=tex-text]{Times New Roman}
%\setmainfont[Mapping=tex-text]{Droid Serif}
%\setmainfont[Mapping=tex-text]{Cambria}
%\setmainfont[Mapping=tex-text]{cm-unicode}
% \setmainfont[Mapping=tex-text]{Gentium}
% \setmainfont[Mapping=tex-text]{GFS Didot}
% \setmainfont[Mapping=tex-text]{Comic Sans MS}
 \setmainfont[Mapping=tex-text]{Ubuntu}
% \setmainfont[Mapping=tex-text]{Myriad Pro}
% \setmainfont[Mapping=tex-text]{CMU Concrete}
%\setmainfont[Mapping=tex-text]{DejaVu Sans}
%\setmainfont[Mapping=tex-text]{KerkisSans}
%\setmainfont[Mapping=tex-text]{KerkisCaligraphic}
% \setmainfont[Mapping=tex-text]{Segoe Print}
% \setmainfont[Mapping=tex-text]{Gabriola}

%------------------------------------------------%

% ------------ Mathematics Fonts ----------------%
%\setmathfont{Asana-Math.ttf}
% -----------------------------------------------%
% Misc
\author{Κασωτάκη Ε. - Σμαραγδάκης Κ.}
\title{www.math24.gr}
% ---------------- Formation --------------------%


\setlength\topmargin{-0.7cm}
\addtolength\topmargin{-\headheight}
\addtolength\topmargin{-\headsep}
\setlength\textheight{26cm}
\setlength\oddsidemargin{-0.54cm}
\setlength\evensidemargin{-0.54cm}
%\setlength\marginparwidth{1.5in}
\setlength\textwidth{17cm}

\RequirePackage[avantgarde]{quotchap}
\renewcommand\chapterheadstartvskip{\vspace*{0\baselineskip}}
\RequirePackage[calcwidth]{titlesec}
\titleformat{\section}[hang]{\bfseries}
{\Large\thesection}{12pt}{\Large}[{\titlerule[0.9pt]}]
%--------------------------------------------------%
%\usepackage{draftwatermark}
%\SetWatermarkText{www.math24.gr}
%\SetWatermarkLightness{0.9}
%\SetWatermarkScale{0.6}
%--------------------------------------------------%


\usepackage{xcolor}
\usepackage{amsthm}
\usepackage{framed}
\usepackage{parskip}

\colorlet{shadecolor}{bluesite!20}

\newtheorem*{orismos}{Ορισμός}

\newenvironment{notation}
  {\begin{shaded}\begin{theorem}}
  {\end{theorem}\end{shaded}}




% Document begins
\begin{document}
\pagestyle{fancy}
\fancyhead{}
\fancyfoot{}
\renewcommand{\headrulewidth}{0pt}
\renewcommand{\footrulewidth}{0pt}


\fancyhead[LO,LE]{
 %\textblockcolor{backsite}
 %\begin{textblock}{5}(-2,-0.55)
  %\rule{0cm}{1cm}
 %\end{textblock}
 \textblockcolor{bluesite}
 \begin{textblock}{5}(-1.5,-0.55)
  \rule{0cm}{1cm}
 \end{textblock}
 %\textblockcolor{bluesite}
 %\begin{textblock}{14}(5.5,-0.55)
 % \rule{0cm}{1cm}
 %\end{textblock}
 \begin{textblock}{0}(-1,-0.25)
 \color{cwhite} \begin{Large}www.math24.gr\end{Large}
 \end{textblock}
\begin{textblock}{0}(16,-0.4)
 \color{cwhite} \includegraphics[height=1.5cm]{math24_logo.png}
 \end{textblock}
\textblockcolor{white}
\begin{textblock}{17}(-1,28)
 \color{backsite} \begin{small}Επιμέλεια: Κασωτάκη Ειρήνη (ikasotaki@gmail.com) - Σχολικό έτος 2020 - 2021\end{small}
 \end{textblock}
}
 \begin{shaded}
\begin{center}
\huge \textbf{Μαθηματικά Α' Γυμνασίου}\\
\end{center} 
\begin{center}
\textbf{"Επαναληπτικές Ασκήσεις"} \hfill \textbf{}
\end{center}
%\subsection*{}
\end{shaded}
%\vspace{2em
%\begin{shaded}
%\begin{center}
%\huge \textbf{Επαναληπτικές Ασκήσεις για το Κεφάλαιο 1}\\
%Πρόσθεση ρητών αριθμών
%\end{center} 
%\textbf{Μαθηματικά Α' Γυμνασίου} \hfill \textbf{Ημερομηνία Παράδοσης : \hspace{2em} }
%\subsection*{Ονοματεπώνυμο :\hfill  \hspace{5em}}
%\end{shaded}
\vspace{2em}
\begin{itemize}
\item 1.3: Δυνάμεις φυσικών αριθμών
\item 1.4: Ευκλείδεια διαίρεση - διαιρετότητα
\item 1.5: Χαρακτήρες διαιρετότητας - ΜΚΔ - ΕΚΠ - Ανάλυση αριθμού σε γινόμενο πρώτων παραγόντων
\item Κεφάλαιο 2: Τα κλάσματα
\item Κεφάλαιο 5: Ποσοστά
\item 6.3: Ανάλογα ποσά - Ιδιότητες ανάλογων ποσών
\item 6.5: Προβλήματα αναλογιών
\item 6.6: Αντιστρόφως ανάλογα ποσά
\item Β.1.3: Μονάδες μήκους – απόσταση σημείων – μέσο ευθύγραμμου τμήματος 
\item Β.1.5: Διχοτόμος γωνίας 
\item Β.1.6: Είδη γωνιών – Κάθετες γωνίες 
\item Β.1.7: Εφεξής – διαδοχικές γωνίες 
\item Β.1.8: Παραπληρωματικές και συμπληρωματικές γωνίες – κατακορυφήν 
\item Β.1.9: Θέσεις ευθειών στο επίπεδο 
\item Β.1.10: Απόσταση σημείου από ευθεία – Απόσταση παραλλήλων 
\item Β.2.3: Μεσοκάθετος ευθύγραμμου τμήματος 
\item Β.2.6: Παράλληλες ευθείες που τέμνονται από μία άλλη ευθεία
\item Β.3.1: Στοιχεία τριγώνου – Είδη τριγώνων 
\item Β.3.2: Άθροισμα γωνιών τριγώνου – Ιδιότητες ισοσκελούς τριγώνου 
\end{itemize}


\section*{Άσκηση 1  \hfill \small{}}
Να κάνετε τις παρακάτω πράξεις:
\begin{enumerate}[1)]
 \item $2^{2}+3^{2}$
 \item $2\cdot 3^{2}$
 \item $3+5^{2}$
 \item $2^{3}\cdot 8 -2\cdot 5$
 \item $(5^{2}-4\cdot 5)+3\cdot5+1^{6}$
 %------------------------------
 \item $(2\cdot 5)^{2}-3^{2}\cdot 3^{2}$
 \item $4+4\cdot 4\cdot2^{2}$
 \item $(1+2)\cdot (3+2^{2})$
 \item $(2+3+4)\cdot (10-2^{3})$
 \item $2\cdot 5^{2}+5\cdot 1^{4}-(10+2-7)$
\end{enumerate}

\section*{Άσκηση 2  \hfill \small{}}
Να βρείτε το $ΕΚΠ$ των παρακάτω αριθμών:
\begin{enumerate}[i)]
 \item $5$ και $20$
 \item $4$ και $3$
 \item $2,3$ και $5$
 \item $2,4$ και $6$
 \item $8,12$ και $15$
\end{enumerate}

\section*{Άσκηση 3  \hfill \small{}}
Να βρείτε το $ΜΚΔ$ των παρακάτω αριθμών:
\begin{enumerate}[i)]
 \item $5$ και $10$
 \item $18$ και $24$
 \item $3$ και $7$
 \item $3,6$ και $9$
 \item $10,20$ και $40$
\end{enumerate}

\section*{Άσκηση 4  \hfill \small{}}
Για τους παρακάτω αριθμούς να εξετάσετε εάν διαιρούνται με τους $2,3,5,10$:
\begin{enumerate}[i)]
 \item $12$
 \item $15$
 \item $21$
 \item $7$
 \item $60$
\end{enumerate}

\section*{Άσκηση 5  \hfill \small{}}
Να εξετάσετε εάν τα παρακάτω κλάσματα είναι ισοδύναμα:
\begin{enumerate}[i)]
 \item $\dfrac{3}{5}$ και $\dfrac{4}{3}$
 \item $\dfrac{2}{6}$ και $\dfrac{1}{3}$
 \item $\dfrac{24}{13}$ και $\dfrac{3}{2}$
 \item $\dfrac{3}{4}$ και $\dfrac{9}{12}$
 \item $\dfrac{1}{5}$ και $\dfrac{3}{15}$
\end{enumerate}

\section*{Άσκηση 6  \hfill \small{}}
Να απλοποιήσετε τα παρακάτω κλάσματα:
\begin{enumerate}[i)]
 \item $\dfrac{5}{15}$
 \item $\dfrac{20}{70}$
 \item $\dfrac{33}{21}$
 \item $\dfrac{18}{48}$
 \item $\dfrac{14}{16}$
\end{enumerate}

\section*{Άσκηση 7  \hfill \small{}}
Να συγκρίνετε τα παρακάτω κλάσματα με τη μονάδα.
\begin{enumerate}[i)]
 \item $\dfrac{6}{5}$
 \item $\dfrac{3}{7}$
 \item $\dfrac{5}{15}$
 \item $\dfrac{20}{20}$
 \item $\dfrac{8}{3}$
\end{enumerate}

\section*{Άσκηση 8  \hfill \small{}}
Να γράψετε τα παρακάτω κλάσματα σε αύξουσα σειρά:
$$ \dfrac{6}{5}, \dfrac{19}{5}, \dfrac{4}{5}, \dfrac{3}{5} $$

\section*{Άσκηση 9  \hfill \small{}}
Να γράψετε τα παρακάτω κλάσματα σε φθίνουσα σειρά:
$$ \dfrac{3}{5}, \dfrac{3}{4}, \dfrac{3}{16}, \dfrac{3}{8} $$

\section*{Άσκηση 10  \hfill \small{}}
Να συγκρίνετε τα παρακάτω κλάσματα:
\begin{enumerate}[i)]
 \item $\dfrac{3}{5}$ και $\dfrac{2}{10}$
 \item $\dfrac{3}{4}$ και $\dfrac{4}{3}$
 \item $\dfrac{1}{5}$ και $\dfrac{2}{3}$
 \item $\dfrac{3}{10}$ και $\dfrac{15}{20}$
 \item $\dfrac{7}{2}$ και $\dfrac{14}{4}$
\end{enumerate}

\section*{Άσκηση 11  \hfill \small{}}
Να κάνετε τις παρακάτω πράξεις:
\begin{enumerate}[i)]
 \item $\dfrac{11}{2}+\dfrac{3}{2}$
 \item $\dfrac{5}{4}+\dfrac{1}{3}$
 \item $\dfrac{4}{2}+\dfrac{5}{3}+\dfrac{1}{4}$
 \item $\dfrac{6}{13}-\dfrac{2}{13}$
 \item $\dfrac{3}{4}-\dfrac{1}{5}$
\end{enumerate}


\section*{Άσκηση 12  \hfill \small{}}
Να κάνετε τις παρακάτω πράξεις:
\begin{enumerate}
 \item $\dfrac{3}{4}\cdot \dfrac{2}{5}$
 \item $\dfrac{1}{2}\cdot \dfrac{3}{2}\cdot \dfrac{2}{3}$
 \item $3\cdot \dfrac{1}{5}$
 \item $4\cdot \dfrac{3}{12}$
 \item $4\cdot \dfrac{1}{5}\cdot \dfrac{2}{3}$
\end{enumerate}

\section*{Άσκηση 13  \hfill \small{}}
Να κάνετε τις παρακάτω πράξεις:
\begin{enumerate}[i)]
 \item $\dfrac{2}{3}+\dfrac{4}{3}\cdot \dfrac{1}{2}$
 \item $(\dfrac{2}{3}+\dfrac{4}{3})\cdot \dfrac{1}{2}$
 \item $\dfrac{2}{3}\cdot (\dfrac{4}{3}+\dfrac{1}{2})$
 \item $\dfrac{1}{3}\cdot (\dfrac{4}{3}-\dfrac{1}{2})$
 \item $\dfrac{2}{3}\cdot (\dfrac{4}{3}-\dfrac{1}{3})$
\end{enumerate}

\section*{Άσκηση 14  \hfill \small{}}
Να βρείτε τους αντίστροφους των παρακάτω αριθμών:
\begin{enumerate}
 \item $\dfrac{5}{3}$
 \item $\dfrac{6}{5}$
 \item $\dfrac{8}{9}$
 \item $\dfrac{1}{7}$
 \item $5$
\end{enumerate}

\section*{Άσκηση 15  \hfill \small{}}
Να κάνετε τις παρακάτω πράξεις:
\begin{enumerate}[i)]
 \item $\dfrac{5}{3}:\dfrac{2}{6}$
 \item $\dfrac{1}{4}:\dfrac{2}{8}$
 \item $(\dfrac{2}{3}:\dfrac{7}{2}):\dfrac{1}{3}$
 \item $(\dfrac{1}{6}:\dfrac{3}{5}):4$
 \item $2:\dfrac{1}{3}$
\end{enumerate}

\section*{Άσκηση 16  \hfill \small{}}
Να μετατρέψετε τα σύνθετα κλάσματα σε απλά:
\begin{enumerate}[i)]
 \item $\dfrac{\dfrac{2}{5}}{\dfrac{4}{7}}$
 \item $\dfrac{\dfrac{1}{2}}{\dfrac{5}{2}}$
 \item $\dfrac{\dfrac{1}{2}+\dfrac{1}{3}}{\dfrac{5}{3}}$
 \item $\dfrac{8}{\dfrac{1}{3}+\dfrac{1}{5}}$
 \item $\dfrac{\dfrac{2}{3}\cdot \dfrac{4}{5}}{7}$
\end{enumerate}

\section*{Άσκηση 17  \hfill \small{}}
Να γράψετε τα παρακάτω κλάσματα ως ποσοστά επί τοις εκατό:
\begin{enumerate}[i)]
 \item $\dfrac{5}{100}$
 \item $\dfrac{4}{10}$
 \item $\dfrac{3}{5}$
 \item $\dfrac{2}{25}$
 \item $\dfrac{7}{20}$
\end{enumerate}

\section*{Άσκηση 18  \hfill \small{}}
Να γράψετε τα παρακάτω  ποσοστά ως κλάσματα:
\begin{enumerate}[i)]
 \item $5\%$
 \item $33\%$
 \item $50\%$
 \item $17\%$
 \item $10\%$
\end{enumerate}

\section*{Άσκηση 19 (εκτός ύλης για το σχολικό έτος 2020 - 2021) \hfill \small{}}
Να συμπληρώσετε τον παρακάτω πίνακα αν γνωρίζετε ότι τα ποσά $x$ και $y$ είναι ανάλογα με συντελεστή αναλογίας $α=4$.
\begin{center}
\begin{tabular}{|c|c|c|c|l|}\hline
\textbf{x} & 1 & 6 &    &  \\\hline
\textbf{y} &   &   & 12 & 16\\\hline
\end{tabular}
\end{center}

\section*{Άσκηση 20 (εκτός ύλης για το σχολικό έτος 2020 - 2021) \hfill \small{}}
Να συμπληρώσετε τον παρακάτω πίνακα αν γνωρίζετε ότι τα ποσά $x$ και $y$ είναι ανάλογα.
\begin{center}
\begin{tabular}{|c|c|c|c|l|}\hline
\textbf{x} & 1 & 2 & 7 & 10 \\\hline
\textbf{y} &   & 4 &   & \\\hline
\end{tabular}
\end{center}

\section*{Άσκηση 21 (εκτός ύλης για το σχολικό έτος 2020 - 2021) \hfill \small{}}
Να εξετάσετε αν τα ποσά $x$ και $y$ που δίνονται στους παρακάτω πίνακες είναι ανάλογα.
\begin{enumerate}[i)]
 \item 
 \begin{center}
\begin{tabular}{|c|c|c|c|}\hline
\textbf{x} & 1 & 2 & 3 \\\hline
\textbf{y} & 3 & 6 & 7 \\\hline
\end{tabular}
\end{center}
 
 \item 
  \begin{center}
\begin{tabular}{|c|c|c|c|}\hline
\textbf{x} & 4 & 6 & 10 \\\hline
\textbf{y} & 6 & 9 & 15 \\\hline
\end{tabular}
\end{center}
\end{enumerate}

\section*{Άσκηση 22 (εκτός ύλης για το σχολικό έτος 2020 - 2021) \hfill \small{}}
Να εξετάσετε αν τα ποσά $x$ και $y$ που δίνονται στους παρακάτω πίνακες είναι αντιστρόφως ανάλογα.
\begin{enumerate}[i)]
 \item 
  \begin{center}
\begin{tabular}{|c|c|c|c|}\hline
\textbf{x} & 2 & 4 & 5 \\\hline
\textbf{y} & 4 & 8 & 3 \\\hline
\end{tabular}
\end{center}
 \item 
  \begin{center}
\begin{tabular}{|c|c|c|c|}\hline
\textbf{x} & 1 & 2 & 3 \\\hline
\textbf{y} & 6 & 3 & 2 \\\hline
\end{tabular}
\end{center}
 \item 
  \begin{center}
\begin{tabular}{|c|c|c|c|}\hline
\textbf{x} & $\dfrac{3}{2}$  & $\dfrac{5}{2}$  & 7 \\\hline
\textbf{y} & $\dfrac{1}{3}$  & $\dfrac{2}{10}$ & $\dfrac{1}{14}$ \\\hline
\end{tabular}
\end{center}
\end{enumerate}

\section*{Άσκηση 23  \hfill \small{}}
Να μετατρέψετε τα παρακάτω μήκη σε μέτρα:
\begin{enumerate}[i)]
 \item $70dm$
 \item $245cm$
 \item $3100mm$
 \item $2,5km$
 \item $50dm$
\end{enumerate}

\section*{Άσκηση 24  \hfill \small{}}
Σε μια ευθεία $ε$ να πάρετε στη σειρά τα σημεία $Α,Β$ και $Γ$ ώστε $ΑΒ=2,5cm$ και $ΒΓ=8,5cm$. Έπειτα να 
βρείτε το μέσο του τμήματος $ΑΒ$.

\section*{Άσκηση 25  \hfill \small{}}
Να σχεδιάσετε μια γωνία $x\hat{O}y=50$\textdegree και να βρείτε τη διχοτόμο της.

\section*{Άσκηση 26  \hfill \small{}}
Να σχεδιάσετε δύο εφεξής και παραπληρωματικές γωνίες.

\section*{Άσκηση 27  \hfill \small{}}
Να σχεδιάσετε δύο εφεξής και συμπληρωματικές γωνίες.

\section*{Άσκηση 28  \hfill \small{}}
Να σχεδιάσετε δύο παράλληλες ευθείες που απέχουν μεταξύ τους $5cm$.

\section*{Άσκηση 29  \hfill \small{}}
Να σχεδιάσετε ένα ευθύγραμμο τμήμα μήκους $8cm$ και στη συνέχεια να σχεδιάσετε τη μεσοκάθετο του.

\section*{Άσκηση 30  \hfill \small{}}
Σε ένα ισοσκελές τρίγωνο η γωνία που είναι απέναντι από τη βάση είναι $96$\textdegree. Να υπολογίσετε τις υπόλοιπες 
γωνίες.

%\section*{Άσκηση 16  \hfill \small{}}
%\begin{multicols}{2}
%\begin{enumerate}[1)]
% \item 
%\end{enumerate}
%\end{multicols}

\end{document}
